%-----------------------------------------------------------------------------------------------------------------------------------------------%
% The MIT License (MIT)
%
% Copyright (c) 2019 Jan Küster
%
% Permission is hereby granted, free of charge, to any person obtaining a copy
% of this software and associated documentation files (the "Software"), to deal
% in the Software without restriction, including without limitation the rights
% to use, copy, modify, merge, publish, distribute, sublicense, and/or sell
% copies of the Software, and to permit persons to whom the Software is
% furnished to do so, subject to the following conditions:
% 
% THE SOFTWARE IS PROVIDED "AS IS", WITHOUT WARRANTY OF ANY KIND, EXPRESS OR
% IMPLIED, INCLUDING BUT NOT LIMITED TO THE WARRANTIES OF MERCHANTABILITY,
% FITNESS FOR A PARTICULAR PURPOSE AND NONINFRINGEMENT. IN NO EVENT SHALL THE
% AUTHORS OR COPYRIGHT HOLDERS BE LIABLE FOR ANY CLAIM, DAMAGES OR OTHER
% LIABILITY, WHETHER IN AN ACTION OF CONTRACT, TORT OR OTHERWISE, ARISING FROM,
% OUT OF OR IN CONNECTION WITH THE SOFTWARE OR THE USE OR OTHER DEALINGS IN
% THE SOFTWARE.
% 
%
%-----------------------------------------------------------------------------------------------------------------------------------------------%




%============================================================================%
%
% DOCUMENT DEFINITION
%
%============================================================================%



%we use article class because we want to fully customize the page and don't use a cv template
\documentclass[10pt]{article} 

\renewcommand{\section}[1]{} % I dont want sections
\renewcommand{\subsection}[1]{} % I dont want subsections

\usepackage{bibentry}
\section{Preamble}
\subsection{packages}

\usepackage{eurosym} % for \euro

\usepackage{fancyhdr} % nom sur le dessus
\usepackage{lastpage} % numéro de page à la fin (référence)

%----------------------------------------------------------------------------------------
% ENCODING
%----------------------------------------------------------------------------------------

% we use utf8 since we want to build from any machine
\usepackage[utf8]{inputenc}  

%----------------------------------------------------------------------------------------
% LOGIC
%----------------------------------------------------------------------------------------

% provides \isempty test
\usepackage{xstring, xifthen}

%----------------------------------------------------------------------------------------
% FONT BASICS
%----------------------------------------------------------------------------------------

% some tex-live fonts - choose your own

%\usepackage[defaultsans]{droidsans}
%\usepackage[default]{comfortaa}
%\usepackage{cmbright}
\usepackage{raleway}      % headings
\usepackage{XCharter}     % body, as in the Sverker proposal
%\usepackage{fetamont}
%\usepackage[default]{gillius}
%\usepackage[light,math]{iwona}
%\usepackage[thin]{roboto} 

% set font default
% body in serif (XCharter)  
\usepackage[T1]{fontenc}

% more font size definitions
\usepackage{moresize}

%----------------------------------------------------------------------------------------
% FONT AWESOME ICONS
%---------------------------------------------------------------------------------------- 

% include the fontawesome icon set
\usepackage{fontawesome5}
% use to vertically center content
% credits to: http://tex.stackexchange.com/questions/7219/how-to-vertically-center-two-images-next-to-each-other
\newcommand{\vcenteredinclude}[1]{\begingroup
\setbox0=\hbox{\includegraphics{#1}}%
\parbox{\wd0}{\box0}\endgroup}

% use to vertically center content
% credits to: http://tex.stackexchange.com/questions/7219/how-to-vertically-center-two-images-next-to-each-other
\newcommand*{\vcenteredhbox}[1]{\begingroup
\setbox0=\hbox{#1}\parbox{\wd0}{\box0}\endgroup}

% icon shortcut
\newcommand{\icon}[3] {        
 \makebox(#2, #2){\textcolor{maincol}{\faIcon{#1}}}
} 

% icon with text shortcut
\newcommand{\icontext}[4]{       
 \vcenteredhbox{\icon{#1}{#2}{#3}}  \hspace{2pt}  \parbox{0.9\mpwidth}{\textcolor{#4}{#3}}
}

% icon with website url
\newcommand{\iconhref}[5]{       
    \vcenteredhbox{\icon{#1}{#2}{#5}}  \hspace{2pt} \href{#4}{\textcolor{#5}{#3}}
}

% icon with email link
\newcommand{\iconemail}[5]{       
    \vcenteredhbox{\icon{#1}{#2}{#5}}  \hspace{2pt} \href{mailto:#4}{\textcolor{#5}{#3}}
}

%----------------------------------------------------------------------------------------
% PAGE LAYOUT  DEFINITIONS
%----------------------------------------------------------------------------------------
\subsection{page layout}
% page outer frames (debug-only)
% \usepackage{showframe}  

% we use paracol to display breakable two columns
\usepackage{paracol}

% define page styles using geometry
\usepackage[a4paper]{geometry}

% remove all possible margins
%\geometry{top=1cm, bottom=1cm, left=1cm, right=1cm}


% remove all possible margins
\geometry{top=1cm, bottom=1.50cm, left=1cm, right=1cm, footskip=1cm}

\usepackage{fancyhdr}
\renewcommand{\headrulewidth}{0pt} 

% space between header and content
% \setlength{\headheight}{0pt}

% indentation is zero
\setlength{\parindent}{0mm}

%----------------------------------------------------------------------------------------
% TABLE /ARRAY DEFINITIONS
%---------------------------------------------------------------------------------------- 

% extended aligning of tabular cells
\usepackage{array}

% custom column right-align with fixed width
% use like p{size} but via x{size}
\newcolumntype{x}[1]{%
>{\raggedleft\hspace{0pt}}p{#1}}%


%----------------------------------------------------------------------------------------
% GRAPHICS DEFINITIONS
%---------------------------------------------------------------------------------------- 

%for header image
\usepackage{graphicx}

% use this for floating figures
% \usepackage{wrapfig}
% \usepackage{float}
% \floatstyle{boxed} 
% \restylefloat{figure}

%for drawing graphics  
\usepackage{tikz}    
\usetikzlibrary{shapes, backgrounds,mindmap, trees}

%----------------------------------------------------------------------------------------
% Color DEFINITIONS
%---------------------------------------------------------------------------------------- 

\subsection{colors}
\usepackage{transparent}
\usepackage{color}
\usepackage{xhfill} % for the line 
% primary color

%% Bleu sobre (photo fin)
\definecolor{maincol}{RGB}{79, 102, 135} 
\definecolor{darkcol}{RGB}{ 54, 67, 73 }
\definecolor{lightcol}{RGB}{236,238,239}

% \definecolor{maincol}{HTML}{001C30}
% \definecolor{darkcol}{HTML}{000000}
% \definecolor{lightcol}{HTML}{E0E0E0}


% \usepackage{pagecolor}
% \definecolor{paper}{HTML}{FEFAF6}

% % dark red
% \definecolor{maincol}{RGB}{118, 36, 49} 
% \definecolor{darkcol}{RGB}{ 54, 67, 73 }
% \definecolor{lightcol}{RGB}{236,238,239}

% %% AMF
% \definecolor{maincol}{RGB}{213, 162, 8}
% \definecolor{darkcol}{RGB}{ 54, 67, 73 } % AMF dark
% \definecolor{lightcol}{RGB}{236,238,239} % AMF pale 

% %% magenta
% \definecolor{maincol}{RGB}{121, 67, 79}
% \definecolor{darkcol}{RGB}{ 54, 67, 73 } 
% \definecolor{lightcol}{RGB}{236,238,239}  

% Package for links, must be the last package used
% hyperref's bibliography patch breaks the bibentry package (its \@lbibitem
% bypasses bibentry's neutralized \item -> "Lonely \item" errors); this CV
% never typesets a bibliography, so opt out of that patch (no-op on old hyperref)
\makeatletter\def\hyper@nopatch@bib{}\makeatother
\usepackage[colorlinks=true,urlcolor=maincol,linkcolor=maincol,citecolor=maincol]{hyperref}

% returns minipage width minus two times \fboxsep
% to keep padding included in width calculations
% can also be used for other boxes / environments
\newcommand{\mpwidth}{\linewidth-\fboxsep-\fboxsep}
 


%============================================================================%
%
% CV COMMANDS
%
%============================================================================%

%----------------------------------------------------------------------------------------
%  CV LIST
%----------------------------------------------------------------------------------------

% renders a standard latex list but abstracts away the environment definition (begin/end)
\newcommand{\cvlist}[1] {
 \begin{itemize}{
 \itemsep0em 
 #1}\end{itemize}
}

%----------------------------------------------------------------------------------------
%  CV TEXT
%----------------------------------------------------------------------------------------

\subsection{CV text}

% base class to wrap any text based stuff here. Renders like a paragraph.
% Allows complex commands to be passed, too.
% param 1: *any
\newcommand{\cvtext}[1] {
 \begin{tabular*}{1\mpwidth}{p{0.98\mpwidth}}
  \parbox{1\mpwidth}{#1}
 \end{tabular*}
}

%----------------------------------------------------------------------------------------
% CV SECTION
%----------------------------------------------------------------------------------------
\subsection{CV section}
% Renders a a CV section headline with a nice underline in main color.
% param 1: section title
\usepackage{needspace} % keeps a section title with its first entry
\newcommand{\cvsection}[1] {
 \needspace{8\baselineskip}
 \vspace{14pt}
 \cvtext{
  {\Large\sffamily\bfseries\textcolor{darkcol}{#1}}\\[-4pt]
  \textcolor{maincol}{ \rule{0.15\textwidth}{2pt}}  \\
    }
    % \xrfill[0ex]{2pt}[lightcol] %to add at 250
}

%----------------------------------------------------------------------------------------
% META SKILL
%----------------------------------------------------------------------------------------

\subsection{CV skill}
% Renders a progress-bar to indicate a certain skill in percent.
% param 1: name of the skill / tech / etc.
% param 2: level (for example in years)
% param 3: percent, values range from 0 to 1
\newcommand{\cvskill}[3] {
 \begin{tabular*}{1\mpwidth}{p{0.53\mpwidth}  r}
   \textcolor{black}{\textbf{#1}} & \textcolor{maincol}{#2}\\
 \end{tabular*}%
 
 \hspace{4pt}
 \begin{tikzpicture}[scale=1,rounded corners=2pt,very thin]
  \fill [lightcol] (0,0) rectangle (1\mpwidth, 0.15);
  \fill [maincol] (0,0) rectangle (#3\mpwidth, 0.15);
   \end{tikzpicture}%
}


%----------------------------------------------------------------------------------------
%  CV EVENT
%----------------------------------------------------------------------------------------
\subsection{CV event}
% Renders a table and a paragraph (cvtext) wrapped in a parbox (to ensure minimum content
% is glued together when a pagebreak appears).
% Additional Information can be passed in text or list form (or other environments).
% the work you did
% param 1: time-frame i.e. Sep 14 - Jan 15 etc.
% param 2:  event name (job position etc.)
% param 3: Customer, Employer, Industry
% param 4: Short description
% param 5: work done (optional)
% param 6: technologies include (optional)
% param 7: achievements (optional)
\newcommand{\cvevent}[7] {
 
 % we wrap this part in a parbox, so title and description are not separated on a pagebreak
 % if you need more control on page breaks, remove the parbox
 \parbox{\mpwidth}{
  \begin{tabular*}{1\mpwidth}{p{0.78\mpwidth}  r}
    \textcolor{black}{\textbf{#2}} & \textcolor{maincol}{#1} \\
   \textcolor{maincol}{\textbf{#3}} & \\
  \end{tabular*}\\[-2pt]
  \ifthenelse{\isempty{#4}}{}{
   \cvtext{#4}\\
  }
 }
 
 \ifthenelse{\isempty{#5}}{}{
  \vspace{3pt}
  \cvtext{#5}
 }
 
 \ifthenelse{\isempty{#6}}{}{
  \vspace{3pt}
  \cvtext{\textbf{Technologies include:}}\\
  \cvtext{#6}
 }
 \ifthenelse{\isempty{#7}}{}{
  \vspace{-6pt}
  \cvtext{\small\textit{Awards:} #7}\\[2pt]
 }
 \vspace{2pt}
}

% audience tag of a talk, one wording for all
\newcommand{\aud}[1]{\textcolor{darkcol!75}{\textit{Audience: #1.}}}

%----------------------------------------------------------------------------------------
%  CV SHORT EVENT: one line, title and time-frame only
\newcommand{\cvshort}[2] {
 \begin{tabular*}{1\mpwidth}{p{0.78\mpwidth}  r}
   \textcolor{black}{\textbf{#2}} & \textcolor{maincol}{#1} \\
 \end{tabular*}\\[1pt]
}

%----------------------------------------------------------------------------------------
%  CV META EVENT
%----------------------------------------------------------------------------------------

\subsection{CV metaevent}
% Renders a CV event on the sidebar
% param 1: title
% param 2: subtitle (optional)
% param 3: customer, employer, etc,. (optional)
% param 4: info text (optional)
\newcommand{\cvmetaevent}[4] {
 \textcolor{maincol} {\cvtext{\textbf{\begin{flushleft}#1\end{flushleft}}}}

 \ifthenelse{\isempty{#2}}{}{
 \textcolor{darkcol} {\cvtext{\textbf{#2}} }
 }

 \ifthenelse{\isempty{#3}}{}{
  \cvtext{{ \textcolor{darkcol} {#3} }}\\
 }

 \cvtext{#4}\\[14pt]
}

%---------------------------------------------------------------------------------------
% QR CODE
%----------------------------------------------------------------------------------------

\subsection{Code QR}
% Renders a qrcode image (centered, relative to the parentwidth)
% param 1: percent width, from 0 to 1
\newcommand{\cvqrcode}[1] {
 \begin{center}
  \includegraphics[width={#1}\mpwidth]{qrcode.jpg}
 \end{center}
}

\section{To do}
%%%%%%%%%%%%%%%%%%%%%%%%%%%% TODO %%%%%%%%%%%%%%%%%%%%%%%%%%%
%
% Add prix coup de coeur enseignement?
% Séparer les communications vulgarisation des communications scientifiques

%============================================================================%
%
%
%
% DOCUMENT CONTENT
%
%
%
%============================================================================%
\begin{document}


\section{Document}
% \pagecolor{paper}
\subsection{biblio}

\nobibliography{biblio}
\bibliographystyle{apalike}

\subsection{Settings}
%%% Footer
\pagestyle{fancy}
\fancyfoot[L]{\small \textcolor{black!35!lightcol}{Olivier Côté, curriculum vitae}}
\fancyfoot[C]{\small \textcolor{black!35!lightcol}{\thepage\ of \pageref*{LastPage}}}
\fancyfoot[R]{\small \textcolor{black!35!lightcol}{Updated \today}}


\columnratio{0.33, 0.33, 0.33} %0.31
\setlength{\columnsep}{2.2em}
\setlength{\columnseprule}{4pt}
\colseprulecolor{lightcol}

%---------------------------------------------------------------------------------------
% TITLE  HEADER
%----------------------------------------------------------------------------------------
\vspace*{-1.25cm} % 1cm of white above the title band, matching the side margins
 
 \section{Header}
{\sffamily\HUGE\textbf{\textcolor{darkcol}{Olivier Côté}}\hspace{8pt}{\LARGE\textcolor{darkcol!70}{ACAS}}}\\[6pt]
{\large\textcolor{darkcol}{PhD Candidate in Actuarial Science, Université Laval}}\\[8pt]
{\small Québec, Canada \enspace\textbar\enspace \href{mailto:olivier.cote.12@ulaval.ca}{olivier.cote.12@ulaval.ca} \enspace\textbar\enspace \href{https://olicoside.github.io/}{olicoside.github.io} \enspace\textbar\enspace \href{https://www.linkedin.com/in/olivier-cote-act/}{linkedin.com/in/olivier-cote-act}}\\[3pt]
{\small Languages: Français (native), English (C1), Svenska (A1) \enspace\textbar\enspace Programming: R, Python, LaTeX}
 


 \vspace*{-0.2cm}
\section{Education}
\cvsection{Education}
% }{}{}{}

\cvevent
{2023--2027}
{PhD in Actuarial Science (GPA 4.17/4.33)}
{Université Laval}
{Fairness in insurance pricing, supervised by Prof. Marie-Pier Côté (Université Laval) and Prof. Arthur Charpentier (UQÀM). Defence expected in summer 2027.}{}{}{}

\cvevent
{2021--2023}
{Master's degree in Actuarial Science, research-based (GPA 4.24/4.33)}
{Université Laval}
{Supervised by Prof. Marie-Pier Côté and Prof. Arthur Charpentier.}{}{}{}

\cvevent
{2018--2021}
{Bachelor's degree in Actuarial Science (GPA 4.17/4.33)}
{Université Laval}
{}{}{}{}

\cvevent
{2023--2026}
{Research visits}
{Stockholm University, University of Toronto, Bayes Business School}
{\cvlist{\item Department of Mathematics, Stockholm University (two weeks, January 2026): invited seminar and research discussions with Prof. Mathias Millberg Lindholm. \item Department of Statistics, University of Toronto (2025): invited seminar and research discussions with Prof. Christopher Blier-Wong. \item Bayes Business School, London (one week, 2023): research discussions with Prof. Andreas Tsanakas.}}{}{}{}







\section{Research Papers}
\cvsection{Research papers}


\cvtext{\bibentry{Annual_rev}}

\vspace*{0.25cm}

\cvtext{\bibentry{Charpentier/al:2025_three}}

\vspace*{0.25cm}

\cvtext{\bibentry{SelectionBias_2024}} 

\vspace*{0.25cm}

\cvtext{\bibentry{A_scalable_toolbox}}

 \vspace*{-0.25cm}

 \cvlist{ \item Winner of the \href{https://actuary.org/ph-d-candidate-olivier-cote-receives-2026-american-academy-of-actuaries-research-award/}{ American Academy of Actuaries Award for Research 2026} \item Selected by the National Association of Insurance Commissioners (NAIC) for its 2026 webinar series for U.S. insurance regulators and their actuarial staff}

 \vspace*{-0.25cm}

\vspace*{0.25cm}

\cvtext{\bibentry{Charpentier/al:2026}} 


\vspace*{0.25cm}

\cvtext{\bibentry{Cote/etal:2024}} 

 \vspace*{-0.25cm}

 \cvlist{ \item Finalist for the \href{https://actuary.org/publication-issue/blog-academy-recognizes-anti-discrimination-paper-with-research-award/}{ American Academy of Actuaries Award for Research 2025}}

 \vspace*{-0.25cm}

\section{Work in progress}
\cvsection{Work in progress}

\cvtext{\textbf{Côté, O.}, Charpentier, A., Côté, M.-P., and Il Idrissi, M. A definition of confounding bias for arbitrary statistical functionals. Working paper, in preparation.}

\vspace*{0.25cm}

\cvtext{\textbf{Côté, O.}, Côté, M.-P., Charpentier, A., and Il Idrissi, M. Profiling actuarial discrimination via causal decomposition. Working paper, presented at the 29th International Congress on Insurance: Mathematics and Economics, Seoul.}



  
\section{Diversity, Equity, and Inclusion Education}
\cvsection{Diversity, equity and inclusion education}

%\vfill\null


 \vspace*{-0.25cm}

 \cvevent
 {2025}
 {Get Comfortable Being Uncomfortable: Dialogue About Favoritism and Fairness}
 {Canadian Statistical Sciences Institute}
 {}{}{}{}

 \cvevent
 {2024}
 {Let's talk about race and collective responsibility}
 {Canadian Statistical Sciences Institute}
 {}{}{}{}

 \cvevent
 {2023}
 {Acting as a sentinel for the prevention of suicide}
 {Centre de prévention du suicide du Québec workshop}
 {}{}{}{}

\cvevent
 {2022}
 {Addressing Conflicts Related to Bias, Privilege, and Identity in the STEM Fields}
 {CANSSI and Academic Impressions workshop}
 {}{}{}{}

%---------------------------------------------------------------------------------------
% WORK EXPERIENCE
%----------------------------------------------------------------------------------------

% \newpage 
\section{Work experiences}
\cvsection{Work experience}


\cvevent
 {2022--2023}
 {Course lecturer, two mandatory undergraduate courses (3 credits each)}
 {Université Laval}
 {\cvlist{\item Linear Models in Actuarial Science (fall 2022, 88 students) \item Statistical Learning in Actuarial Science (winter 2023, 54 students)}}{}{}{Faculty of Science and Engineering Teaching Star, one per course (fall 2022 and winter 2023).}
%Students left over 93\% positive feedbacks about course planning, pedagogy, teacher and evaluation.

\cvevent
 {2019--2025}
 {American football coach}
 {Falcons of Cégep de Lévis}
 {American football coach for different teams for more than ten years (2012--2025), including seven years in collegiate football at the Cégep de Lévis (2019--2025): 
 \cvlist{\item Tight ends and Fullbacks coach (2019--2025), responsible for the development of 10 athletes \item Supervisor for a bachelor's degree intern in sports intervention (2023) \item Volunteer mathematics tutor for student-athletes
 }}{}{}{RSEQ Coaching Staff of the Year 2023 and 2024, Sportsmanship Award 2022 (D1) and 2023 (D2).}

% \item Assistant offensive line coach (2019--2021)
%  \item Assistant offensive coordinator (2022--present)
 \cvevent
 {2019--2024}
 {Teaching assistant}
 {Université Laval}
 {Statistical Learning in Actuarial Science (5 times), Probabilistic Analysis of Actuarial Risks (3 times), Linear Models in Actuarial Science (once), internship evaluation (once).}{}{}{}
 


\cvevent
 {2019--2021}
 {Actuarial intern}
 {Desjardins General Insurance Group}
 {Four internships: property ratemaking (2019), business analysis (2020 and 2021) and research (2021).}{}{}{}

\section{Supervision and mentoring}
\cvsection{Supervision and mentoring}

\cvevent
 {2025--2026}
 {Mentor of a master's student, University of Texas}
 {Hickman Scholars Program network}
 {Advised an essay on an insurance regulatory issue and connected the student with a Texas regulator.}{}{}{}

\cvevent
 {2023}
 {Co-supervisor of a master's student, Université Laval}
 {With Prof. Marie-Pier Côté}
 {One-year project on fairness in insurance. The student now works as an actuary.}{}{}{}

\cvevent
 {Ongoing}
 {Supervisor of a PhD student, Université Laval}
 {Mitacs Accelerate partnership}
 {Helped design the thesis project of the student who continues the industrial partnership.}{}{}{}

\section{Engagement with practice and policy}
\cvsection{Engagement with practice and policy}

\cvevent
 {2025--2026}
 {Online course on fair insurance pricing, English and French versions}
 {Casualty Actuarial Society (in post-production)}
 {Commissioned by the CAS. Project lead and script writer, coordinating four parties: a research assistant, the CAS task force on biases in insurance pricing, the CAS educational development team and a production company. Two hours of two-minute videos, interactive activities and quizzes, co-presented with M.-P. Côté and A. Charpentier.}{}{}{}

\cvevent
 {2023--2026}
 {Regulators, insurers and public bodies}
 {Canada, United States, Switzerland}
 {\cvlist{\item Met insurance regulators in person in Canada, the United States and Switzerland. \item Provided input, in person, to the Swiss Federal Statistical Office and the ETH Zurich AI Policy Lab on a governmental framework for fair predictive models in employment. \item Presented at the internal annual actuarial conferences of three insurers. \item Three companies contacted me for technical details on implementing the toolbox with their own data.}}{}{}{}
% Four internships in different teams for Desjardins General Insurance Group. I developed business knowledge, communication skills and data science expertise : 


\section{Refereeing}
\cvsection{Refereeing}

\cvshort{2026--present}{\textit{European Actuarial Journal}}
\cvshort{2026--present}{\textit{Annals of Actuarial Science}}
\cvshort{2024--present}{\textit{Journal of Risk and Insurance}}




\section{Scholarships and Awards}
\cvsection{Scholarships and awards}



\cvevent
 {2026}
 {Award for Research}
 {American Academy of Actuaries}
 {Winner (honorarium US\$7~500, $\approx$ 75~000 SEK) for \textit{A scalable toolbox for exposing indirect discrimination in insurance rates}. Annual international competition for papers by early-career scholars.}{}{}{}

\cvevent
 {2026}
 {Travel Grant}
 {NSERC General Fund}
 {Travel grant to support conference participation and dissemination of research (\$1~300 CAD, $\approx$ 9~100 SEK).}{}{}{}

\vspace*{-0.05cm}

\cvevent
 {2026}
 {Actuarial Research Conference 2026 Travel Grant}
 {American Academy of Actuaries}
 {Travel grant to support conference participation and dissemination of research (US\$1~500, $\approx$ 15~000 SEK).}{}{}{}

\cvevent
 {2026}
 {Student travel grant}
 {\href{https://crdm.ulaval.ca/2025/10/09/bourse-de-mobilite-etudiante-du-crdm/}{Big Data Research Center}}
 {Travel grant to support conference participation and dissemination of research (\$2~000 CAD, $\approx$ 14~000 SEK).}{}{}{}

\cvevent
 {2025}
 {Canada Graduate Research Scholarship -- Doctoral}
 {Natural Sciences and Engineering Research Council of Canada}
 {National competition, \$120~000 CAD ($\approx$ 841~000 SEK) over three years, held from May 2025.}{}{}{}

\cvevent
 {2024, 2025, 2026}
 {Hickman Scholars Program}
 {Society of Actuaries}
 {International competition, US\$20~000 ($\approx$ 199~000 SEK) per year for three years (US\$60~000 in total, $\approx$ 598~000 SEK). Every annual renewal is competitive.}{}{}{}


\cvevent
 {2024}
 {Doctoral Training Scholarships}
 {Fonds de recherche du Québec -- Nature et technologies}
 {Provincial competition, \$100~000 CAD ($\approx$ 701~000 SEK) offered over four years, held from May 2024 to May 2025, when the NSERC doctoral scholarship began.}{}{}{}



\cvevent
 {2024}
 {Mitacs Accelerate}
 {Mitacs}
 {Research partnership with a large Canadian P\&C insurer, \$84~000 CAD ($\approx$ 589~000 SEK) over three years (2024--2027), in which I am the lead student researcher. Selected after a scientific peer review and an internal Mitacs review.}{}{}{}




 \cvevent
 {2023}
 {Doctorate attraction grant Florent Toureille}
 {Chaire d'actuariat de l'Université Laval}
 {Institutional competition based on leadership and excellence, \$30~000 CAD ($\approx$ 210~000 SEK) over three years.}{}{}{}


\cvevent
 {2023}
 {Leadership and Engagement Scholarship}
 {Université Laval}
 {Institutional competition in the sports sector, \$30~000 CAD ($\approx$ 210~000 SEK) over three years.}{}{}{}



\cvevent
 {2023}
 {Charles Brindamour Graduate Scholarship for Excellence}
 {Université Laval, funded by Charles Brindamour (CEO of Intact Financial Corporation)}
 {Institutional competition based on excellence, one award a year, \$10~000 CAD ($\approx$ 70~000 SEK).}{}{}{}



\cvevent
 {2022}
 {François-Laviolette Scholarship}
 {Université Laval}
 {Institutional competition for responsible artificial intelligence, \$7~500 CAD ($\approx$ 53~000 SEK).}{}{}{}
 


\cvevent
 {2021}
 {Grant of attraction to the Master's}
 {Chaire d'actuariat de l'Université Laval}
 {Institutional competition recognizing research potential, \$15~000 CAD ($\approx$ 105~000 SEK).}{}{}{}



\cvevent
 {2021}
 {CAS Trust Scholarship}
 {Casualty Actuarial Society}
 {International competition, US\$2~500 ($\approx$ 25~000 SEK).}{}{}{}



\cvevent
 {2021}
 {Alexander Graham Bell Canada Graduate Scholarship -- Master's}
 {Natural Sciences and Engineering Research Council of Canada}
 {National competition, \$17~500 CAD ($\approx$ 123~000 SEK).}{}{}{}



\cvevent
 {2018}
 {Lieutenant Governor's Youth Medal}
 {Cégep de Lévis}
 {Awarded for outstanding involvement and academic results, on nomination by the teaching staff.}{}{}{}



%
%



%---------------------------------------------------------------------------------------
% EDUCATION
%----------------------------------------------------------------------------------------

\section{Professional designations}
\cvsection{Professional designations}

%\vfill\null

\cvevent
 {2025}
 {Associate of the Casualty Actuarial Society (ACAS)}
 {Casualty Actuarial Society}
 {}{}{}{}

\cvevent
 {2025}
 {Professionalism Workshop}
 {Canadian Institute of Actuaries}
 {}{}{}{}

\cvevent
 {2020--2025}
 {MAS-I, MAS-II, Exam 5, Exam 6C, Online Courses 1 and 2}
 {Casualty Actuarial Society}
 {}{}{}{}
 

 \cvevent
 {2019--2020}
 {Exam P, FM and IFM }
 {Society of Actuaries and Canadian Institute of Actuaries (credits)}
 {}{}{}{}
% \vfill\null
%
%\vfill


\newpage

\section{Academic Communications}
\cvsection{Academic communications}


\cvevent{2026}{Operationalizing Fairness in Actuarial Pricing: From Principle to Practice}{Casualty Actuarial Society webinar}{Sole presenter for a 90-minute webinar to 950 registrants. \aud{practitioners}}{}{}{}

\vspace*{-0.05cm}

\cvevent{2026}{The Trilemma of Fairness in P\&C Pricing}{61st Actuarial Research Conference (AAA invited session)}{Invited presentation. \aud{decision-makers}}{}{}{}

\vspace*{-0.05cm}

\cvevent{2026}{Profiling Actuarial Discrimination via Causal Decomposition}{29th International Congress on Insurance: Mathematics and Economics, Sungkyunkwan University, Seoul (South Korea)}{Academic presentation. \aud{academics}}{}{}{}

\vspace*{-0.05cm}

\cvevent{2026}{Le trilemme de l'équité en tarification IARD}{Association des actuaires en IARD du Québec (AAIARD)}{Presentation to P\&C actuaries. \aud{practitioners}}{}{}{}

\vspace*{-0.05cm}

\cvevent{2026}{Profiling Actuarial Discrimination via Causal Decomposition}{STATQAM seminar, Université du Québec à Montréal}{Invited seminar covering causal decomposition and confounding bias. \aud{academics}}{}{}{}

\vspace*{-0.05cm}

\cvevent{2026}{Profiling Actuarial Discrimination via Causal Decomposition}{Centre de recherche en données massives, Université Laval}{Research talk. \aud{academics}}{}{}{}

\vspace*{-0.05cm}

\cvevent{2026}{Profiling Actuarial Discrimination via causal decomposition}{Actuarial and Financial Mathematics Conference, Brussels (Belgium)}{Poster presenter. \aud{academics}}{}{}{}

\vspace*{-0.05cm}

\cvevent
 {2026}
 {Bias in Assessing Financial Risk: Insights from the 2025 Research Award}
{\href{https://www.actuary.org/event/bias-assessing-financial-risk-insights-2025-research-award}{American Academy of Actuaries}}
 {Invited panelist as finalist for the 2025 Award for Research on Bias in Assessing Financial Risk. Webinar to 445 actuaries and practitioners. \aud{practitioners}}{}{}{}


\vspace*{-0.05cm}

\cvevent
 {2026}
 {A Causal Perspective on Fairness in Insurance Pricing}
 {Department of Mathematics, Stockholm University}
 {Invited seminar talk during research stay. \aud{academics}}{}{}{}



\vspace*{-0.05cm}


\cvevent{2026}{A Scalable Toolbox for Exposing Indirect Discrimination In Insurance Rates}{ASTIN Bulletin Conference, Zurich (Switzerland)}{Speaker at one of the fairness sessions of the first ASTIN Bulletin Conference. \aud{academics}}{}{}{}


\vspace*{-0.05cm}

\cvevent
 {2025}
 {Operationalizing Fairness in Actuarial Pricing: From Principle to Practice}
 {Department of Statistics, University of Toronto}
 {Invited seminar during research visit. \aud{academics}}{}{}{}


\vspace*{-0.05cm}

\cvevent
 {2025}
 {Operationalizing Fairness in Actuarial Pricing: From Principle to Practice}
 {Casualty Actuarial Society Annual Meeting, Austin (United States) \& Canada Connect Day, Toronto (Canada)}
 {Invited presenter to over 150 (Austin) and over 100 (Toronto) actuaries and policymakers. \aud{practitioners}}{}{}{}


\vspace*{-0.05cm}

\cvevent
 {2025}
 {Two invited presentations at the 60th Actuarial Research Conference}
 {American Academy of Actuaries panel \& Casualty Actuarial Society session, Toronto (Canada)}
 {Guest panelist and presenter (with Marie-Pier and Arthur) to audiences from actuarial organizations, academia, industry, and policymakers. \aud{decision-makers}}{}{}{}



\vspace*{-0.05cm}

\cvevent{2025}{A Scalable Toolbox for Exposing Indirect Discrimination in Insurance Rates}{Perspectives on Actuarial Risks in Talks of Young Researchers (PARTY), Liverpool (U.K.)}{Presenter to 60 doctoral students from 29 countries. \aud{academics}}{}{}{}


\vspace*{-0.05cm}

\cvevent
 {2025}
 {Une boîte à outils pour dévoiler la discrimination indirecte dans les tarifs d'assurance}
 {Chaire ACTIONS \& Milliman R\&D, Paris (France)}
 {Guest co-presenter (with Arthur) to audiences from academia and industry (80 participants at the Chaire ACTIONS). \aud{academics and practitioners}}{}{}{}


\vspace*{-0.05cm}

\cvevent{2024}{Perspective causale sur les discriminations dans les modèles prédictifs en assurance}{Institut intelligence et données}{Guest presenter to over 100 people from academia. \aud{academics}}{}{}{}


\vspace*{-0.05cm}

\cvevent{2024}{A Fair price to pay: exploiting causal graphs for fairness in insurance}{27th International Congress on Insurance: Mathematics and Economics}{Academic presentation in Chicago, IL, United States. \aud{academics}}{}{}{}


\vspace*{-0.05cm}

\cvevent{2024}{A Fair price to pay: exploiting causal graphs for fairness in insurance}{Workshop on fairness and discrimination in insurance}{Guest presenter to over 60 people. \aud{academics, practitioners and regulators}}{}{}{}


\vspace*{-0.05cm}

\cvevent{2024}{Perspective causale sur la discrimination indirecte en assurance générale}{Autorité des marchés financiers}{Presentation to decision-makers at Québec's financial authority, followed by interviews with ten regulators on their needs. \aud{decision-makers}}{}{}{}






\vspace*{-0.05cm}

\cvevent{2023}{Équité et discrimination en assurance générales}{Intact Financial Corporation}{Presentation to CEO, CLO and CDO of Intact Financial Corporation. \aud{decision-makers}}{}{}{}


\vspace*{-0.05cm}

\cvevent{2023}{Causal inference and Fairness in Insurance Pricing}{Insurance Data Science Conference}{Presenter for \href{https://insurancedatascience.org/}{Insurance Data Science Conference \faExternalLink*} in London (UK). \aud{academics}}{}{}{}











\section{Public Outreach}
\cvsection{Public outreach}

\cvevent{2023--2026}{L'avenir à travers une boule de cristal mathématique : Le travail d'un actuaire}{Cégep de Lévis}{Guest presenter for CEGEP students. \aud{students}}{}{}{}

\cvevent{2024}{Pour une intelligence artificielle équitable en assurance}{Université Laval News}{I collaborated for a  \href{https://nouvelles.ulaval.ca/2024/05/17/pour-une-intelligence-artificielle-equitable-en-assurance-a:0fa1aae9-47cd-478a-a61e-4ea4d441533f}{scientific communication article \faExternalLink*}. \aud{general public}}{}{}{}

\cvevent{2024}{Perspective causale sur la discrimination indirecte en assurance générale}{Institut Intelligence et données}{4th research competition on data intelligence. $\approx$ 50 graduate students. \aud{general public}}{}{}{}

\cvevent{2024}{Perspective causale sur la discrimination indirecte en assurance générale}{Institut de valorisation des données}{\href{https://www.youtube.com/watch?v=MSz1W8Ri_y4&list=PLnI094tgrXpwdY1U0eYSuRv3-IzdtnG0C&index=7}{I presented my research \faExternalLink*} during IVADO Digital Future Day 2024. Over 250 people. \aud{general public}}{}{}{}






\cvevent{2023}{L'IA pour calculer le prix de votre assurance auto: poétique ou pas éthique?}{Big Data Research Center}{Presentation during the event \textit{lightning talks 180}. \aud{general public}}{}{}{Public's choice award}



\cvevent{2023}{Artificial intelligence to prevent unfair discrimination in auto insurance}{IVADO}{I authored an 800-word  \href{https://ivado.ca/en/research-projects/prevent-unfair-discrimination-in-auto-insurance/}{scientific communication article \faExternalLink*}. \aud{general public}}{}{}{}








% \section{Other academic contributions}
% \cvsection{Other academic contributions}




% ajoute rle lien pedagogic material 

%---------------------------------------------------------------------------------------
% Commitment/engagement 
%----------------------------------------------------------------------------------------

\vspace*{-0.35cm}

\section{Organizing}
\cvsection{Organizing}

 \cvevent
 {2026}
 {Workshop on fairness and discrimination in insurance}
 {Université Laval, 6--7 October 2026, 90 participants}
 {Member of the organizing committee, host of the second day, panel moderator and presenter.}{}{}{}

 \cvevent
 {2025}
 {Quantact Student Day}
 {Province-wide event, 3 October 2025}
 {Organizer: recruited a five-member committee, coordinated seven presenters and hosted sixty students.}{}{}{}

\vspace*{-0.35cm}

\section{Committees}
\cvsection{Committees}

 \cvevent
 {2025--present}
 {Institut Intelligence et Données, Université Laval}
 {Voting member of the board}
 {}{}{}{}

 \cvevent
 {2022--present}
 {Actuarial Science Graduate Program Committee of Université Laval}
 {Voting member of the Committee}
 {}{}{}{}
 \cvevent
 {2022--2025}
 {Big Data Research Center of Université Laval}
 {Student voting member of the Committee}
 {}{}{}{}

 \cvevent
 {2021--2024}
 {Actuarial Science Student Investment Fund of Université Laval}
 {Voting member of the committee}
 {}{}{}{}

 

%---------------------------------------------------------------------------------------
% Research afiliation
%----------------------------------------------------------------------------------------


\vspace*{-0.35cm}

\section{Research affiliations}
\cvsection{Research affiliations}
\cvevent
 {2021--present}
 {Institut Intelligence et Données, Big Data Research Center}
 {Université Laval}
 {}{}{}{}

\cvevent
 {2021--present}
 {Quantact research laboratory, Centre de recherches mathématiques}
 {Montréal}
 {}{}{}{}


%---------------------------------------------------------------------------------------
% Examens Professionnels
%----------------------------------------------------------------------------------------


\vspace*{-0.35cm}

\section{Other}
\cvsection{Other}

 \cvevent{2025}{Roundtable on raising awareness about youth bullying}{\href{https://www.youtube.com/watch?v=PsvMgb9Yje8}{Mammouth hebdo - Téléquébec \faExternalLink* }}{}{}{}{}

 \cvevent{2022--2023}{Actuarial School representative at the open house}{Faculty of Science and Engineering}{\cvlist{\item Winter 2023 ( $\approx$ 30 visits from students accompanied by their parents) \item Fall 2022 ( $\approx$ 50 visits from students accompanied by their parents) }}{}{}{}

% \item Clean and jerk : 275 pounds for 1 rep (2023), 245 for 3 reps (2023) \item Running : half marathon in 109 minutes (2024), 10K in 46 minutes (2025), 5K in 22 minutes (2025)}}{}{}


% \item Offensive lineman of the year for Pointe-Lévy high school (2012) \item All-star team as an offensive lineman in Cadet D2 (2012)

\end{document}

